% !TeX root = ../main.tex
\section{Datasets}
\label{sec:datasets}

\subsection{What makes a dataset useful here}

For transcribing a guitar part from a commercial song, four properties decide
how useful a dataset is: (i) whether the audio is \emph{electric} guitar with
realistic amplification, (ii) whether labels include the \emph{string} (not just
the pitch) and the \emph{techniques}, (iii) whether the guitar appears inside a
\emph{mix}, and (iv) the \emph{licence}, which decides whether a model trained
on it can be distributed. No public dataset has all four.

A recurring term is \term{DI} (direct input): the clean signal straight from the
guitar's pickup, before any amplifier. A DI recording can be played through any
number of amplifier simulations afterwards (\term{re-amping} or
\term{rendering}), which turns one labelled performance into many labelled
tones. This is how EGDB, GOAT and our own training data multiply their size.

\subsection{The datasets}

\paragraph{Audio with string labels.}
GuitarSet~\cite{xi2018guitarset} records acoustic guitar with a hexaphonic
pickup: per-string notes and pitch contours for 360 excerpts by six players
(00--05), each excerpt played as accompaniment (``comp'') and as a melody
(``solo''). It is the standard benchmark, and its small size makes overfitting a
recurring concern~\cite{zang2024synthtab}.
EGDB~\cite{chen2022egdb} provides 240 tablatures played on electric guitar with
onset, pitch and string labels, as DI and through five amplifier simulations,
plus a few real recordings; EGDB-PG~\cite{chen2025egdbpg} re-renders it through
many amplifier-cabinet configurations.
GOAT~\cite{loth2025goat} has 5.9\,h of DI electric guitar with Guitar Pro tabs
that include techniques, and 29.5\,h of amplifier renders.
EGSet12~\cite{pedroza2024tones} has 12 real performances through a miked
amplifier with string labels.
IDMT-SMT-Guitar~\cite{kehling2014idmt} has isolated notes and licks labelled
with expression styles (bend, slide, vibrato, harmonics, dead notes).

\paragraph{Audio with note labels only.}
GAPS~\cite{riley2024gaps} aligns published classical guitar scores to 14\,h of
recordings by more than 200 performers: the largest real guitar set with note
labels. Guitar-TECHS~\cite{pedroza2025guitartechs} records electric guitar at
once as DI, miked amplifier and room microphones, with MIDI labels from a
hexaphonic MIDI pickup, covering techniques, chords, scales and musical
excerpts.

\paragraph{Symbolic and synthetic.}
DadaGP~\cite{sarmento2021dadagp} is 26\,181 Guitar Pro songs (no audio), the
main resource for learning fingering from human tabs; it is available for
research on request. SynthTab~\cite{zang2024synthtab} renders DadaGP through
commercial guitar plugins. GuitarFlow~\cite{loth2025guitarflow} renders tabs
with a sampler and then transforms the result toward realistic electric guitar.

\paragraph{Mixes.}
MUSDB18~\cite{rafii2017musdb18} (150 songs) and
MoisesDB~\cite{pereira2023moisesdb} (240 songs, with guitar sub-stems split into
acoustic, clean electric and distorted electric) are separation datasets: real
multitrack songs, no note labels. MuScriptor~\cite{rouard2026muscriptor} built a
private set of 170\,000 real recordings (over 11\,000\,h) with notes aligned by
audio-to-score synchronisation~\cite{muller2021sync}.

\begin{table}[h]
\centering
\caption{Datasets and how this project uses them. Licences as recorded in
\file{scripts/train/README.md} and \file{docs/transcription_roadmap.md}; NC =
non-commercial, ND = no derivatives.}
\label{tab:datasets}
\small
\begin{tabularx}{\textwidth}{@{}>{\raggedright}p{2.3cm}>{\raggedright}p{3.0cm}ccc>{\raggedright}p{2.0cm}X@{}}
\toprule
Dataset & Audio & String & Tech. & Mix & Licence & Our use \\
\midrule
GuitarSet & acoustic, mic + hexaphonic & yes & contours & no & CC BY 4.0 & Players 00--02 train; 03--05 test (never fitted on) \\
EGDB & electric DI + 5 amp sims & yes & no & no & none stated & Train: ids not multiple of 5; test: multiples of 5 and the real clips \\
EGSet12 & real electric, miked amp & yes & no & no & CC BY 4.0 & Test only \\
GOAT & electric DI + amp renders & yes (tabs) & yes & no & NC, on request & Official test split; never in a distributed model \\
Guitar-TECHS & electric DI, amp, mics & no & yes & no & CC BY 4.0 & Chords/scales train; excerpts and techniques test \\
IDMT-SMT-Guitar & electric notes and licks & yes & yes & no & CC BY-NC-ND & Test only (techniques) \\
GAPS & classical, real & no & no & no & -- & Pre-training of the GAPS model (by its authors) \\
DadaGP & none (symbolic) & yes & yes & -- & research, on request & \planned{} fingering pre-training \\
MUSDB18 & real stems & no & no & yes & NC research & Backing for synthetic band mixes \\
MoisesDB & real stems & no & no & yes & check terms & Real-mix benchmark (Section~\ref{sec:bench-moises}) \\
\bottomrule
\end{tabularx}
\end{table}

\begin{inapp}[: where the data lives]
None of the datasets is in the repository. They live under
\file{~/.cache/gva_test} (about 60\,GB), prepared by the scripts in
\file{scripts/train/} (\code{prepare_data.py}, \code{make_mixes.py},
\code{prepare_moisesdb.py}, \code{align_labels.py}), which write manifests
(\code{.tsv} files listing audio and notes files) that the Rust benchmark
\code{eval_manifest} and the Python benchmarks read.
\end{inapp}

\begin{caution}[: the missing dataset]
There is no public set of commercial band recordings with reference tabs, which
is exactly the use case. MuScriptor showed how one could be built (align
existing tabs to recordings, reject bad alignments). The app has the tool for
this, \code{scripts/train/align_tabs.py} \done, but no aligned and hand-checked
song pairs yet (Section~\ref{sec:bench-real}).
\end{caution}
